%%% Local Variables:
%%% mode: latex
%%% TeX-master: "../doc"
%%% coding: utf-8
%%% End:
% !TEX TS-program = pdflatexmk
% !TEX encoding = UTF-8 Unicode
% !TEX root = ../doc.tex
% !TEX root = ../doc.tex
\section{Discussion of Research Questions}

The empirical results presented in Chapter~\ref{ch:results} and the prototype described in Chapter~\ref{ch:prototype} together provide a basis for revisiting the five research questions posed in Section~\ref{sec:research_questions_and_thesis}. This section addresses each question in turn, drawing on the quantitative trust outcomes, the helpfulness rankings, and the qualitative reflections provided by participants.

\subsubsection*{RQ1: Does the provision of explainability information help users to appropriately trust or distrust the output of an LLM?}

The results provide a clear affirmative answer, though with an important qualification: explainability information does change user trust, but not uniformly, and not always in the direction one might initially consider desirable.

Across all XAI conditions, trust was lower than the plain LLM baseline (Mdn\,=\,4.0), which produced the highest and most stable trust ratings of any condition. At first glance, this might suggest that XAI methods are counterproductive. However, this interpretation is too hasty. The baseline condition presented participants with fluent, confident outputs and no further information, precisely the conditions under which miscalibrated over-trust is most likely to develop. As discussed in Section~\ref{sec:trust_in_context_of_llms}, LLMs are susceptible to hallucination and dishonesty, and their outputs are delivered with the same confident fluency regardless of accuracy. Users who default to high trust in the absence of any explanatory signal are not necessarily better calibrated; they are simply more credulous.

Viewed through this lens, the reduction in trust produced by XAI methods is not a failure as long as that the trust reduction is directed toward outputs that genuinely warrant scepticism. The qualitative reflection data supports this reading: among the 18.9\% of participants who reported a change in trust after the survey, a substantial proportion described a shift toward conditional or more nuanced trust rather than blanket scepticism (Table~\ref{tb:example_responses_trust}). Responses in the \textbf{Conditional Trust} and \textbf{Greater Understanding} categories indicate that participants were not simply disillusioned by the survey, but emerged with a more differentiated view of when and how much to rely on LLM outputs.

This pattern aligns with the theoretical framework of appropriate calibration in human-centred XAI: an explanation that reduces inflated trust in an unreliable output improves calibration even if it does not increase overall trust. The central empirical finding of this study is therefore that XAI methods move user trust in the right direction: away from unconditional reliance, even when the resulting trust levels appear lower in absolute terms \cite{zhang_effect_2020}.

\subsubsection*{RQ2: How do different XAI methods differentially affect lay users' trust levels?}

The three methods produced meaningfully distinct trust profiles, confirming that the choice of explanation paradigm matters for lay users.

Chain-of-Thought prompting together with Confidence Scores produced the lowest median trust of all conditions (Mdn\,=\,3.0), with mean ratings across aspects ranging from 2.7 to 3.5. This sharp decline from the baseline was consistent across trust aspects and visible at the individual trajectory level (Figure~\ref{fig:individual}), where participants who entered the CoT condition with high baseline trust showed the steepest drops. At the same time, CoT received the highest number of first-place helpfulness rankings ($n = 45$), making it simultaneously the most preferred and the least trust-inducing method. This paradox is interpretively coherent: CoT exposes the model's reasoning process in a form users find informative and engaging, and this very transparency reveals inconsistencies or errors that appropriately reduce trust. The method therefore functions as a high-information, trust-calibrating explanation type, consistent with its theoretical characterization in Section 5.1.

Confidence indication produced a median trust score identical to CoT (Mdn\,=\,3.0) but received the second-highest helpfulness ranking ($n = 43$ first-place votes) and the most second-place votes across all methods. Unlike CoT, confidence scores do not expose reasoning content; they communicate uncertainty directly through a numerical or colour-coded signal. The similar trust outcome with a lower information load suggests that confidence scores function as a more efficient but shallower calibration mechanism. Users can quickly identify outputs the model is uncertain about without having to evaluate a reasoning chain. This makes confidence indication particularly well-suited to high-throughput, time-constrained use cases.

SHAP-based attribution and counterfactual reasoning produced intermediate trust scores (Mdn\, = \,3.4) but the lowest helpfulness rankings by a clear margin ($n = 27$ first-place votes, most third-place votes). This combination of moderate trust with low perceived usefulness is consistent with prior literature indicating that feature-importance visualizations are poorly understood by non-technical audiences \cite{chromik_i_2021, kaur_interpreting_2020}. For the predominantly non-technical sample in this study, SHAP values appear to have been legible enough to modestly reduce trust relative to the baseline, but not comprehensible enough to be experienced as genuinely helpful. The technical framing of token attribution, even when accompanied by a plain-language counter-analysis, did not translate into a subjectively useful explanation for most participants.

The ``How an LLM works'' explanation (EXPL) occupied an intermediate position (Mdn\,=\,3.6) and outperformed SHAP and CONF on most trust aspects, particularly perceived competence and reliability. This suggests that providing users with a general conceptual model of LLM behaviour, even without reference to any specific output, can itself modulate trust in a constructive direction.

\subsubsection*{RQ3: Is XAI a reliable indicator of whether the model is actually functioning correctly, or can it create misleading impressions of transparency?}

The answer depends critically on which XAI method is considered, and the evidence from both the study and the literature reviewed in Chapter~\ref{ch:prototype} warrants caution.

Confidence scores are the most straightforwardly reliable indicator available within the methods studied here. Derived directly from token-level log-probability distributions, they reflect a genuine property of the model's output distribution rather than a post-hoc narrative. However, calibration is not guaranteed: as discussed in Section~\ref{sec:selected_xai_methods}, aligning expressed confidence with actual response quality is a non-trivial problem \cite{tao_when_2024}, and a model may assign high probability to a fluently generated but factually incorrect token. Confidence scores are therefore a useful but imperfect signal.

SHAP-based attribution, while grounded in cooperative game theory and methodologically principled, introduces a different reliability concern. The KernelExplainer used in this study computes an approximation of Shapley values rather than exact values, and the approximation quality is controlled by the number of masking samples. This can be a trade-off that, under computational constraints, may reduce the precision of the attributions. More fundamentally, SHAP explains which parts of the input influenced the model's output given its current parametrization, but it does not indicate whether that influence was appropriate or whether the resulting output is correct. A prompt token may carry a high SHAP value precisely because the model has learned a spurious association.

Chain-of-Thought prompting raises the most serious faithfulness concerns of all three methods. As detailed in Section~\ref{sec:faithfulness_problem},  there exists substantial variation in CoT faithfulness across tasks, with larger models showing \textbf{less} faithful reasoning than smaller ones because they can arrive at confident answers without relying on the generated trace\cite{lanham_measuring_2023}. CoT reasoning can be systematically biased by irrelevant input features while the stated reasoning provides a different and misleading account\cite{turpin_language_2023}. The practical implication is that a CoT explanation may create a compelling but inaccurate impression of the model's reasoning process. This is a concern that is particularly acute in high-stakes domains where users are most likely to seek an explanation in the first place.

Taken together, these findings suggest that no current XAI method for LLMs provides a fully reliable indicator of correct functioning for lay users. Each method introduces its own failure modes, and users should be encouraged to treat XAI outputs as useful but imperfect signals rather than authoritative accounts of model behaviour. This limitation reinforces the importance of the EXPL condition: equipping users with a conceptual understanding of how LLMs work may be at least as important as providing per-output explanations, because it helps users interpret those explanations with appropriate scepticism.

\subsubsection*{RQ4: How can general users be made more aware of the appropriate degree of trust to place in AI-generated information?}

The findings suggest two complementary pathways, which together constitute the practical contribution of this thesis.

The first is the combined CoT and confidence prototype described in Chapter~\ref{ch:prototype}. By pairing the high-information, trust-calibrating properties of Chain-of-Thought reasoning with the lightweight, immediately interpretable confidence signal, the prototype targets the central tension identified in the results. Users can follow the model's reasoning at a step-by-step level while receiving a colour-coded confidence indicator for each step that allows them to identify where the model is most uncertain without requiring any technical knowledge of model architecture. The automatic regeneration feature further reduces the cognitive burden on users: steps below a configurable confidence threshold are regenerated without manual intervention, directing user attention toward outputs that are more likely to be reliable. This design operationalizes the concept of appropriate calibration in a form accessible to lay users.

The second pathway is educational: providing users with a plain-language account of how LLMs work. The EXPL condition, which delivered a short video or text explanation of LLM mechanics, produced consistently higher trust scores than SHAP and CONF and recovered more of the ground lost relative to the baseline. This suggests that conceptual framing (helping users understand that LLMs are statistical pattern-matchers rather than reasoning agents) has a meaningful effect on trust calibration that is complementary to per-output explanation methods. Practical deployments of LLM-based tools would therefore benefit from incorporating on-boarding material that conveys this understanding, alongside real-time explanatory features.

\subsubsection*{RQ5: What is the computational and usability overhead of add-on XAI methods, and does this overhead limit their practical deployment?}

The three methods differ substantially in their computational profiles, and this difference has direct implications for practical deployment.

Confidence indication carries the lowest overhead of the three methods. Token-level log-probabi\-li\-ties are a byproduct of the standard auto-regressive generation process and can be extracted at negligible additional cost by enabling the \localtt{output\_scores} flag in the inference API. No additional model calls are required, and the aggregate statistics can be computed in microseconds. The only usability overhead is the design of the confidence display, which must communicate uncertainty in a form that is immediately interpretable without requiring technical knowledge.

Chain-of-Thought prompting requires a modified system prompt and a parsing step to extract structured reasoning steps from the model's output, but it does not require any additional model inference calls. The main computational cost is the additional output tokens generated by the reasoning trace, which extends generation time roughly in proportion to the length of the reasoning chain. In the survey application, the JSON-based extraction proved brittle and required a repair utility; the prototype redesigned this as a delimiter-based parser, which is both more reliable and faster. The main usability overhead is the interface design required to render reasoning steps in a form that is scannable and comprehensible for non-expert users.

SHAP-based attribution carries by far the highest computational overhead of the three methods. The KernelExplainer requires $N$ masked forward passes over the prompt, where $N$ is a user-configurable sample count that controls the approximation quality. In the survey application, each SHAP inference request required substantially longer processing time than the CoT or confidence endpoints, and the resulting bar charts were generated server-side with an additional Matplotlib call. A second inference call was then required to produce the plain-language counter-analysis. This accumulation of computational costs makes SHAP attribution impractical for real-time, interactive use cases without significant optimization, which could for example be achieved through caching, batching or the use of lighter surrogate approximations. Combined with its low perceived helpfulness among lay users, these costs make SHAP the least viable of the three methods for consumer-facing deployment in its current form.

\section{Outlook}

Several directions for future work emerge from the findings of this study.

The most immediate extension concerns the breadth of XAI methods evaluated. This study examined three methods chosen to cover distinct explanation paradigms, but the XAI landscape for LLMs is considerably wider. Future work could include attention-based visualizations, narrative gamification frameworks \cite{you_gamifying_2024}, or counterfactual-only conditions without the accompanying SHAP scores, allowing their individual contributions to trust calibration to be assessed. The EXPL condition, which performed well in this study, could itself be elaborated into a more interactive on-boarding experience rather than a static video or text.

A second direction concerns the faithfulness problem identified for Chain-of-Thought prompting. The prototype introduced per-step confidence scores as a complementary signal, but this does not address the underlying issue that the stated reasoning may not reflect the model's actual computation. Approaches such as SelfIE \cite{chen_selfie_2024}, which grounds explanations in the model's internal embeddings rather than its output tokens, represent a technically more faithful alternative. A key challenge for future research is to develop interfaces that make such embedding-level explanations accessible to non-technical users. That means bridging the gap between mechanistic faithfulness and lay usability that currently separates these two strands of XAI research.

Third, the sample in this study, while demographically broad, was drawn from a panel of registered survey participants and skewed toward administrative, healthcare, and retired demographics. Future studies should evaluate XAI methods with more technically literate populations to assess whether the low helpfulness of SHAP is truly an artefact of the non-technical sample or a more general finding. Within-subject designs with domain-matched tasks (e.g. presenting legal queries to legal professionals and medical queries to clinicians) would allow the interaction between domain expertise and explanation type to be examined more precisely.

Fourth, the prototype developed in this thesis demonstrated the CoT and confidence combination on a locally hosted open-source model. Future work could evaluate whether the same combination is effective when applied to commercial black-box APIs, where token-level probabilities may not be directly accessible. In such settings, calibrated verbal hedging (the model expressing its own uncertainty in natural language) represents a practical alternative whose effect on trust calibration has not yet been rigorously evaluated in a within-subject design.

Finally, the regulatory context established by the EU AI Act creates a practical imperative for XAI research that is oriented toward demonstrable compliance rather than technical benchmarks. Future work should examine how explanation methods can be designed and evaluated in ways that satisfy transparency requirements not only for technically sophisticated auditors but for the non-specialist users who are most directly affected by AI-assisted decisions.

\section{Summary}

This thesis set out to investigate whether and how Explainable AI methods can help lay users calibrate their trust in LLM outputs. The empirical user study with $n = 175$ prior LLM users found that all three XAI methods (Chain-of-Thought prompting, SHAP-based attribution with counterfactual reasoning, and confidence indication) reduced trust relative to a plain LLM baseline, but that this reduction reflects appropriate recalibration rather than failure. Chain-of-Thought emerged as the most preferred and most trust-disrupting method, confidence indication as an efficient and moderately preferred alternative, and SHAP as technically grounded but poorly suited to non-technical audiences. A general explanation of how LLMs work produced competitive trust outcomes, underscoring the value of conceptual framing alongside per-output explanation.

Building on these findings, a prototype was developed that combines Chain-of-Thought reasoning with per-step confidence scores and automatic regeneration of low-confidence steps. This combination targets the central design gap identified in the literature and the study results: explanation methods that are technically grounded enough to be reliable, yet rendered in forms accessible enough to be useful for general users. Together, the empirical findings and the prototype contribute a set of concrete design recommendations for LLM-based applications in which appropriate user trust is not merely desirable but, in high-stakes domains, essential.